% !TEX root = ../main.tex
% =====================================================================
%  CHƯƠNG 2 — CƠ SỞ LÝ THUYẾT VÀ LỰA CHỌN LINH KIỆN
% ---------------------------------------------------------------------
%  Người phụ trách : [TODO: Họ tên - MSSV]
%  Tiền tố label   : c2      (fig:c2:..., tab:c2:..., eq:c2:..., sec:c2:...)
%  Thư mục ảnh     : image/chuong2/
%  File TLTK       : references/chuong2.bib   (key: c2_...)
%  Độ dài gợi ý    : 6–10 trang
%  Mục tiêu        : Trình bày lý thuyết nền và LÝ DO lựa chọn từng linh kiện
%                    (so sánh, bảng thông số trích từ datasheet, có trích dẫn).
% =====================================================================
\chapter{Cơ sở lý thuyết và lựa chọn linh kiện}
\label{chap:c2}

% ---------------------------------------------------------------------
\section{Tổng quan về Load cell}
\label{sec:c2:loadcell}

\subsection{Cấu tạo}
\label{sec:c2:loadcell-cau-tao}
% Gợi ý: thân đàn hồi, strain gauge, vỏ bọc. Có hình minh họa.

\TODO{Trình bày cấu tạo Load cell.}

\subsection{Nguyên lý hoạt động}
\label{sec:c2:loadcell-nguyen-ly}
% Gợi ý: cầu Wheatstone, quan hệ điện áp ra theo khối lượng, có phương trình đánh số.

Điện áp ngõ ra của Load cell tỉ lệ tuyến tính với khối lượng đặt lên cảm biến:
\begin{equation}
    \Vout = \Vin \cdot RO \cdot \frac{m}{m_{\max}}
    \label{eq:c2:vout-loadcell}
\end{equation}
trong đó $RO$ là độ nhạy định mức (\unit{\milli\volt\per\volt}), $m$ là khối lượng cần đo
và $m_{\max}$ là tải trọng định mức của Load cell.

\TODO{Giải thích nguyên lý cầu Wheatstone, kèm hình.}

% ---------------------------------------------------------------------
\section{Lựa chọn Load cell}
\label{sec:c2:chon-loadcell}
% Gợi ý: so sánh 2--3 loại có bán tại VN, chọn 1 loại, bảng thông số từ datasheet.

\TODO{Lý do lựa chọn Load cell.}

\begin{table}[H]
    \centering
    \caption{Thông số kỹ thuật Load cell \TODO{model}}
    \label{tab:c2:thong-so-loadcell}
    \renewcommand{\arraystretch}{1.3}
    \begin{tabular}{|l|c|}
        \hline
        \textbf{Thông số} & \textbf{Giá trị} \\ \hline
        Tải trọng định mức           & \qty{1}{\kilo\gram} \\ \hline
        Độ nhạy định mức (RO)        & \TODO{} \unit{\milli\volt\per\volt} \\ \hline
        Trở kháng ngõ vào            & \TODO{} \unit{\ohm} \\ \hline
        Trở kháng ngõ ra             & \TODO{} \unit{\ohm} \\ \hline
        Điện áp hoạt động            & \TODO{} \unit{\volt} \\ \hline
    \end{tabular}
\end{table}

% ---------------------------------------------------------------------
\section{Lựa chọn bộ khuếch đại}
\label{sec:c2:chon-opamp}
% Gợi ý: INA (instrumentation amplifier) cho tín hiệu vi sai mV;
% OPAMP chính xác cho các khối V-I, đệm, cộng. Nêu thông số quan trọng:
% Vos, drift, CMRR, GBW, nguồn cấp. Trích dẫn datasheet bằng \cite{}.

Bộ khuếch đại dụng cụ được lựa chọn dựa trên các thông số trong datasheet~\cite{c2_ina333}.

\TODO{So sánh và chọn INA / OPAMP, kèm bảng thông số.}

% ---------------------------------------------------------------------
\section{Cảm biến nhiệt độ cho mạch bù nhiệt}
\label{sec:c2:cam-bien-nhiet}

\TODO{Giới thiệu RTD/NTC/IC cảm biến nhiệt và lý do lựa chọn.}

% ---------------------------------------------------------------------
\section{Chuẩn tín hiệu dòng 4--20mA}
\label{sec:c2:chuan-4-20ma}

\TODO{Ưu điểm của chuẩn 4--20mA, khái niệm ngõ ra thả nổi.}

% ---------------------------------------------------------------------
\section{Lựa chọn ADC}
\label{sec:c2:chon-adc}
% Gợi ý: số bit tối thiểu từ độ phân giải yêu cầu: 2^N >= tầm đo / độ phân giải.

\TODO{Tính số bit tối thiểu và chọn ADC / phương án ADC.}

% ---------------------------------------------------------------------
\section{Tổng hợp linh kiện sử dụng}
\label{sec:c2:tong-hop}

\begin{table}[H]
    \centering
    \caption{Tổng hợp các linh kiện chính được lựa chọn}
    \label{tab:c2:tong-hop-linh-kien}
    \renewcommand{\arraystretch}{1.3}
    \begin{tabularx}{\textwidth}{|c|l|l|X|}
        \hline
        \textbf{STT} & \textbf{Khối} & \textbf{Linh kiện} & \textbf{Lý do lựa chọn} \\ \hline
        1 & Cảm biến        & \TODO{} & \TODO{} \\ \hline
        2 & Khuếch đại      & \TODO{} & \TODO{} \\ \hline
        3 & Chuyển đổi V-I  & \TODO{} & \TODO{} \\ \hline
        4 & ADC             & \TODO{} & \TODO{} \\ \hline
    \end{tabularx}
\end{table}
